% Source: https://artofproblemsolving.com/community/c3h2526790_escalator_problem
% Stable identity: aops-topic-2526790-q01
Relative to the ground, the escalator carries each person down at $100$ steps per minute.

\textbf{Maria} moves up at $152-100=52$ steps per minute, so she climbs the $50$ steps in
\[\dfrac{50}{52}=\dfrac{25}{26}\text{ min}=\dfrac{25}{26}\cdot60\text{ s}=\dfrac{750}{13}\text{ s}\approx57.7\text{ s},\]
and leaves at the top.

\textbf{Dee} moves at $13-100=-87$ steps per minute, that is, he is carried \emph{down} at $87$ steps per minute. He starts at the bottom, so he cannot get any higher: he is carried straight back off the escalator at the bottom, at time $0$.

So Dee leaves first, at the bottom, and Maria leaves at the top
\[\boxed{\dfrac{750}{13}\approx57.7\text{ seconds}}\]
later.
